\documentclass[11pt]{article}

\usepackage[T1]{fontenc}
\usepackage[utf8]{inputenc}
\usepackage{lmodern}
\usepackage{amsmath,amssymb,amsthm,mathtools}
\usepackage{booktabs}
\usepackage{enumitem}
\usepackage[a4paper,margin=1in]{geometry}
\usepackage{microtype}
\usepackage{xcolor}
\usepackage[hidelinks]{hyperref}

\newtheorem{theorem}{Theorem}[section]
\newtheorem{proposition}[theorem]{Proposition}
\newtheorem{lemma}[theorem]{Lemma}
\newtheorem{corollary}[theorem]{Corollary}
\theoremstyle{definition}
\newtheorem{definition}[theorem]{Definition}
\newtheorem{problem}[theorem]{Open problem}
\theoremstyle{remark}
\newtheorem{remark}[theorem]{Remark}

\newcommand{\Sn}{\mathcal S_n}
\newcommand{\E}{\mathbb E}
\newcommand{\R}{\mathbb R}
\newcommand{\Cut}{\operatorname{Cut}}
\newcommand{\sgn}{\operatorname{sign}}
\newcommand{\SDP}{\operatorname{SDP}}
\newcommand{\statusproved}{\textcolor{green!40!black}{\textsc{proved}}}
\newcommand{\statuscomputed}{\textcolor{blue!65!black}{\textsc{exhaustive computation}}}
\newcommand{\statusnumerical}{\textcolor{orange!80!black}{\textsc{numerical}}}
\newcommand{\statusopen}{\textcolor{red!70!black}{\textsc{open}}}
\newcommand{\statusrefuted}{\textcolor{purple!70!black}{\textsc{refuted}}}

\title{\textbf{Extremal Seidel Quadratic Forms on the Boolean Cube}\\
\large Bounds, Exact Values, Barriers, and Open Problems}
\author{Marcus Miccelli}
\date{3 August 2026}

\hypersetup{
  pdftitle={Extremal Seidel Quadratic Forms on the Boolean Cube},
  pdfauthor={Marcus Miccelli},
  pdfsubject={Bounds, exact values, barriers, and open problems},
  pdfkeywords={Seidel matrix, Boolean cube, quadratic form, discrepancy}
}

\begin{document}
\maketitle

\begin{abstract}
For $n\ge2$, define
\[
 F(n)=\min_{a_{ij}\in\{-1,1\}}
 \max_{x\in\{-1,1\}^n}
 \left|\sum_{1\le i<j\le n}a_{ij}x_ix_j\right|,
 \qquad L_n=F(n)n^{-3/2}.
\]
The existence and value of $\lim_n L_n$ remain open.  This report records a
 set of internally audited bounds, structural identities, exact computations,
 counterexamples, and continuation problems.  The main universal result is
 \[
  0.336493364432\ldots\le\liminf_{n\to\infty}L_n
  \le\limsup_{n\to\infty}L_n\le\frac12.
 \]
 The lower endpoint follows from a field-plus-spin rounding, independently
 reconstructed here after its appearance in an external AI-assisted research
 ledger; the upper endpoint follows from Paley conference matrices and
 principal restriction.  The earlier squared-Gram argument remains as an
 explicit finite-order bound.  Exact values are reported through order $15$.
 A second part of the report
maps the remaining absolute-field problem: it is proved in several diffuse,
bounded-component, bounded-degree, and twin-fibre regimes, while a genuinely
dense overlapping-correlation regime remains.  The report also records
barriers that rule out several natural spectral, compositional,
superadditive, and low-degree Fourier approaches.  For composition, it is
enough to prove a power-saving defect on comparable block sizes; explicit
rectangular estimates already give sublinear defects when one block is
$o(n)$ and power savings when it is polynomially smaller.
\end{abstract}

\noindent\fbox{\parbox{0.95\textwidth}{%
\textbf{Status.} This is an independent research report, not a journal-refereed
paper.  Statements are labelled as proved, exhaustive computation, numerical,
open, or refuted.  ``Internally audited'' means that normalizations and stated
 dependencies were re-derived or checked independently during the project; it
 does not replace external peer review.  The public repository reproduces its
 own exact search only through order $6$.  A pinned external package for order
 $11$ was rerun locally, as were external witnesses and solver cases through
 order $14$; the remaining larger results retain their historical evidence
 labels.  External source code is cited rather than vendored.
}}

\tableofcontents

\section{Problem, normalization, and interpretations}

\begin{definition}
A \emph{Seidel matrix} of order $n$ is a matrix $A\in\R^{n\times n}$ such that
$A=A^{\mathsf T}$, $A_{ii}=0$, and $A_{ij}\in\{-1,1\}$ for $i\ne j$.
The set of such matrices is denoted by $\Sn$.
\end{definition}

For $A\in\Sn$ and $x\in\{-1,1\}^n$, put
\[
 Q_A(x)=\sum_{i<j}A_{ij}x_ix_j=\frac12x^{\mathsf T}Ax,
 \qquad
 M(A)=\max_x|x^{\mathsf T}Ax|.
\]
Also write
\[
 M_+(A)=\max_xx^{\mathsf T}Ax,
 \qquad
 M_-(A)=-\min_xx^{\mathsf T}Ax.
\]
The normalization ledger is therefore
\begin{equation}\label{eq:normalization}
 F(n)=\frac12\min_{A\in\Sn}M(A),
 \qquad L_n=\frac{F(n)}{n^{3/2}}.
\end{equation}
Every factor of two in this report is with respect to
\eqref{eq:normalization}.

The problem has several equivalent or neighbouring interpretations.
It is a discrepancy problem for a two-colouring of $E(K_n)$, a deterministic
Ising ground-state minimization, and the minimum $L^\infty$ norm of a
 homogeneous order-two Rademacher chaos.  The asymptotic question was posed by
 Ivanisvili on MathOverflow \cite{ivanisvili}.  Volberg's Boolean Sidon quantity
$T_{n,2}$ is exactly $F(n)$ \cite{volberg}; Astashkin and Lykov establish the
$n^{3/2}$ order up to universal constants in a more general weighted setting
\cite{astashkin-lykov}.  The induced-subset discrepancy studied by
Erd\H{o}s and Spencer is related but not identical \cite{erdos-spencer}.

Switching $A\mapsto DAD$, where $D$ is diagonal with entries $\pm1$, and
simultaneous row/column permutation both preserve $M$, $M_\pm$, and the cut
quantity below.  Thus the natural finite objects are switching classes
(two-graphs), modulo permutation when classification is desired.

\section{Principal conclusions}

Let $\phi$ and $\Phi$ be the standard normal density and distribution
function.  Let $t_*>0$ maximize
$4\phi(t)(2\Phi(t)-1)$, equivalently
\[
 2\phi(t_*)=t_*\bigl(2\Phi(t_*)-1\bigr),
 \qquad t_*=0.8769009855\ldots,
\]
and put
\[
 \gamma_*:=2\phi(t_*)\bigl(2\Phi(t_*)-1\bigr)
 =0.336493364432\ldots .
\]

\begin{theorem}[universal window]\label{thm:window}
For every $n\ge2$ and every $A\in\Sn$,
\begin{equation}\label{eq:squared-gram-bound}
 M(A)\ge \frac{2}{\pi}n(n-1)
 \arcsin\!\left(\frac1{\sqrt{n-1}}\right)
 \ge \frac{2}{\pi}n\sqrt{n-1}.
\end{equation}
In addition, every sequence $A_n\in\mathcal S_n$ satisfies
\[
 \liminf_{n\to\infty}\frac{M(A_n)}{n^{3/2}}\ge2\gamma_*.
\]
There are Seidel matrices of every sufficiently large order, after principal
restriction from nearby Paley conference orders, for which
\[
 F(n)\le\left(\frac12+o(1)\right)n^{3/2}.
\]
Hence
\[
 \gamma_*\le\liminf L_n\le\limsup L_n\le\frac12.
\]
\end{theorem}

The existence of the limit is not asserted.  The exact data do not presently
distinguish convergence from persistent arithmetic oscillation.

\subsection{Exact values}

 The following values are classified as \statuscomputed.  Orders $2$ through
 $6$ are regenerated by the public code.  Order $11$ now has a separately
 published standalone enumeration that was rerun during this audit; public
 external witness and solver reruns also corroborate orders $12$ through $14$.
 The remaining entries above order $6$ retain the historical exhaustive or
 threshold-emptiness evidence described in Section~\ref{sec:computation}.

\begin{center}
\begin{tabular}{c|rrrrrrrrrrrrrr}
\toprule
$n$&2&3&4&5&6&7&8&9&10&11&12&13&14&15\\
\midrule
$\min_A M(A)$&2&6&8&8&10&18&20&24&26&34&36&40&42&54\\
$F(n)$&1&3&4&4&5&9&10&12&13&17&18&20&21&27\\
\bottomrule
\end{tabular}
\end{center}

The minimizer is unique up to switching and permutation at orders $12$, $13$,
and $14$: respectively a doubling of $C_6$, a vertex deletion of $C_{14}$,
and the Paley conference matrix $C_{14}$.  At order $15$, the all-positive
one-vertex extension of $C_{14}$ is optimal with $M=54$.  Monotonicity and
congruence imply $\min_A M(A)\ge56$ at order $16$.

For the cut discrepancy introduced below, the exact minimum for
$n=3,\ldots,14$ is
\[
 2,4,4,5,8,10,12,13,16,18,20,21.
\]
At order $15$ it is either $24$ or $26$; the exhaustive route needed to
separate the two exceeded the available memory.

\section{Core identities and lower bounds}

\begin{lemma}[fill-in]\label{lem:fill-in}
For $A\in\Sn$ and $u\in\{0,\pm1\}^n$,
$|u^{\mathsf T}Au|\le M(A)$.
\end{lemma}

\begin{proof}
Complete the zero coordinates of $u$ by independent Rademacher signs to
obtain $x\in\{\pm1\}^n$.  All mixed and off-support terms have expectation
zero; no diagonal term survives because $A_{ii}=0$.  Thus
$u^{\mathsf T}Au=\E[x^{\mathsf T}Ax]$, whose absolute value is at most
$M(A)$.
\end{proof}

The zero diagonal is essential.  For
$\left(\begin{smallmatrix}4&-1\\-1&-3\end{smallmatrix}\right)$, the Boolean
maximum is $3$, but the vector $(1,0)$ has quadratic value $4$.

For a partition $[n]=S\sqcup T$, define
\[
 \Cut(A)=\max_{S\sqcup T=[n]}
 \max_{u\in\{\pm1\}^S,\,v\in\{\pm1\}^T}
 u^{\mathsf T}A_{S,T}v.
\]

\begin{theorem}[cut identity]\label{thm:cut-identity}
For every $A\in\Sn$,
\[
 M_+(A)+M_-(A)=4\Cut(A).
\]
In particular,
$\max_x|Q_A(x)|\ge\Cut(A)$.
\end{theorem}

\begin{proof}
For independently chosen $x,y\in\{\pm1\}^n$, put
$u=(x+y)/2$ and $w=(x-y)/2$.  Their supports partition $[n]$ and
\[
 x^{\mathsf T}Ax-y^{\mathsf T}Ay=4u^{\mathsf T}Aw.
\]
Conversely every signed cut arises from such a pair.  Maximizing $x$ and $y$
independently gives the identity.
\end{proof}

\begin{proposition}[fixed-split first moment]\label{prop:first-moment}
For every $1\le s<n$,
\[
 \Cut(A)\ge(n-s)\E|S_s|,
\]
where $S_s$ is a sum of $s$ independent Rademacher variables.
\end{proposition}

\begin{proof}
Fix $|S|=s$.  For $u\in\{\pm1\}^S$, optimize the signs on $T$ to obtain
$\|A_{S,T}^{\mathsf T}u\|_1$, then average over $u$.  Each of the $n-s$
coordinates is distributed as $S_s$.
\end{proof}

For odd $s=2k+1$,
\[
 \E|S_s|=(2k+1)4^{-k}\binom{2k}{k}
 \ge\sqrt{\frac2\pi}\sqrt{s}.
\]
The parity restriction matters: the displayed inequality is false for even
$s$.  A sharpened Wallis inequality due to Kazarinoff \cite{kazarinoff},
together with a six-residue finite check, absorbs the rounding loss in choosing
an odd integer near $n/3$.  This gives the finite theorem
\begin{equation}\label{eq:first-lower}
 F(n)\ge \max_{1\le s<n}(n-s)\E|S_s|
 \ge \frac{2\sqrt2}{3\sqrt{3\pi}}n^{3/2}
 =0.3071059106\ldots\,n^{3/2}.
\end{equation}
The stronger asymptotic bound is Theorem~\ref{thm:window}.

\subsection{Two-orientation squared-Gram rounding}

Put $q=n-1$ and $B=A/\sqrt q$.  Since $B_{ii}=0$ and
$(B^2)_{ii}=1$, the matrices
\[
 X_\sigma=\frac12(I+\sigma B)^2,
 \qquad \sigma\in\{-1,1\},
\]
are positive semidefinite correlation matrices.  Hyperplane-round each
$X_\sigma$.  For $i\ne j$, write
$b=B_{ij}$ and $c=(B^2)_{ij}/2$.  The sign-correlation identity yields
terms involving $\arcsin(c+\sigma b)$.  Averaging the two orientations uses
\[
 b\bigl[\arcsin(c+b)-\arcsin(c-b)\bigr]
 \ge2|b|\arcsin|b|.
\]
Indeed, the even density $(1-t^2)^{-1/2}$ has least mass on a fixed-length
interval when that interval is centered at zero.  Summing the
 $n(n-1)$ identical values $|b|=q^{-1/2}$ proves
 \eqref{eq:squared-gram-bound}.

\subsection{Field-plus-spin improvement}

The rounding scheme in this subsection appeared in the independent research
ledger \cite{beaumont-ledger}.  The proof below was reconstructed from the
definitions, with the normalization, pair expansion, replacement error, and
final factor of two checked afresh.

\begin{lemma}[spectral bootstrap]\label{lem:spectral-bootstrap}
For every $A\in\Sn$,
\[
 M(A)\ge\frac{\|A\|_2^3}{n}.
\]
\end{lemma}

\begin{proof}
Let $Av=\lambda v$, where $\|v\|_2=1$ and
$|\lambda|=\|A\|_2$.  The eigenvector equation gives
\[
 |\lambda|\|v\|_\infty\le\|v\|_1\le\sqrt n.
\]
Put $c=\|v\|_\infty^{-1}$ and choose independent signs $X_i$ with
$\E X_i=cv_i$.  The zero diagonal makes all variance terms disappear, so
\[
 \E X^{\mathsf T}AX=c^2v^{\mathsf T}Av.
\]
Taking absolute values and using
$c\ge\|A\|_2/\sqrt n$ proves the claim.
\end{proof}

\begin{proof}[Proof of the field-plus-spin part of
Theorem~\ref{thm:window}]
It is enough to consider a sequence for which $M(A)=O(n^{3/2})$; otherwise
the claimed lower bound is automatic along the relevant subsequence.  Put
$m=n-1$ and $\lambda=\|A\|_2$.  Lemma~\ref{lem:spectral-bootstrap} gives
$\lambda=O(n^{5/6})$.  For $i\ne j$, set
\[
 r_{ij}=\frac{A_{ij}(A^2)_{ij}}m.
\]
Since $\operatorname{tr}A^2=nm$,
\begin{equation}\label{eq:mean-decorrelation}
 \frac1{nm}\sum_{i\ne j}r_{ij}^2
 =\frac{\|A^2-mI\|_F^2}{nm^3}
 \le\frac{\operatorname{tr}A^4}{nm^3}
 \le\frac{\lambda^2}{m^2}=O(n^{-1/3}).
\end{equation}
Thus the normalized correlations $r_{ij}$ tend to zero on average in
absolute value.

Fix $t>0$ and $\tau>0$.  Let $\xi_1,\ldots,\xi_n$ be independent
Rademacher signs, let $Z_1,\ldots,Z_n$ be independent standard Gaussians,
and for $\sigma\in\{-1,1\}$ define
\begin{equation}\label{eq:fieldspin-witness}
 Y_i^\sigma=\sgn\!\left(
  \sigma\frac{(A\xi)_i}{\sqrt m}+t\xi_i+\tau Z_i
 \right).
\end{equation}
Write
\[
 \psi_\tau(s)=2\Phi(s/\tau)-1,
\]
and introduce the bounded smooth functions
\[
 \alpha(u)=\frac{\psi_\tau'(u+t)+\psi_\tau'(u-t)}2,
 \qquad
 \beta(v)=\frac{\psi_\tau(v+t)-\psi_\tau(v-t)}2.
\]

We record the pair calculation.  Fix $i\ne j$ and switch the matrix so that
$A_{ij}=1$ and the common coefficients in row $i$ are all $1$.  If the
corresponding row-$j$ coefficients are $c_k\in\{-1,1\}$, put
\[
 U=\frac1{\sqrt m}\sum_{k\ne i,j}\xi_k,
 \qquad
 V=\frac1{\sqrt m}\sum_{k\ne i,j}c_k\xi_k,
 \qquad
 d=\sum_{k\ne i,j}c_k.
\]
Here $d=A_{ij}(A^2)_{ij}$ before switching.  After integrating out the
Gaussian dithers in \eqref{eq:fieldspin-witness}, Taylor expansion in
$\varepsilon=m^{-1/2}$ shows that the two orientations have equal zeroth
order terms and
\begin{equation}\label{eq:pair-response}
 A_{ij}\left(
  \E Y_i^+Y_j^+-\E Y_i^-Y_j^-
 \right)
 =4\varepsilon\E[\alpha(U)\beta(V)]+O_{t,\tau}(\varepsilon^2).
\end{equation}
The factor $4$ comes from the two endpoint spins and the two opposite
orientations.  All derivatives used in the remainder are bounded because
$\tau$ is fixed.

The pair $(U,V)$ is a sum of $m-1$ independent two-dimensional Rademacher
vectors, has coordinate variances $1-1/m$, and covariance $d/m=r_{ij}$.
A two-dimensional Lindeberg replacement therefore has uniform error
$O_{t,\tau}(m^{-1/2})$: each summand has size $O(m^{-1/2})$, there are
$O(m)$ summands, and third derivatives are bounded.  Adjusting the variances
to one costs $O_{t,\tau}(m^{-1})$.  Gaussian covariance interpolation then
gives, uniformly even near singular covariance,
\begin{equation}\label{eq:pair-replacement}
 \E[\alpha(U)\beta(V)]
 =\overline\alpha\,\overline\beta
  +O_{t,\tau}\!\left(m^{-1/2}+|r_{ij}|\right),
\end{equation}
where, for $G\sim N(0,1)$ and
$\phi_v(s)=(2\pi v)^{-1/2}e^{-s^2/(2v)}$,
\[
 \overline\alpha=\E\alpha(G)=2\phi_{1+\tau^2}(t),
 \qquad
 \overline\beta=\E\beta(G)
 =2\Phi\!\left(\frac{t}{\sqrt{1+\tau^2}}\right)-1.
\]

Average $\sigma$ uniformly in \eqref{eq:fieldspin-witness}.  For every realization,
$\sigma(Y^\sigma)^{\mathsf T}AY^\sigma\le M(A)$.  Summing
\eqref{eq:pair-response}--\eqref{eq:pair-replacement} over the $nm$ ordered
pairs, and using \eqref{eq:mean-decorrelation} to make the average of
$|r_{ij}|$ vanish, yields
\[
 \frac{M(A)}{n\sqrt m}
 \ge
 4\phi_{1+\tau^2}(t)
 \left[
  2\Phi\!\left(\frac{t}{\sqrt{1+\tau^2}}\right)-1
 \right]-o(1).
\]
First let $n\to\infty$, then let $\tau\downarrow0$, and finally optimize
over $t>0$.  This proves
\[
 \liminf_n\frac{M(A_n)}{n^{3/2}}
 \ge4\phi(t_*)\bigl(2\Phi(t_*)-1\bigr)=2\gamma_*.
\]
Choosing a minimizing $A_n$ and applying
\eqref{eq:normalization} gives the asserted lower bound for $L_n$.
\end{proof}

\subsection{Spectral-edge rounding}

Let
$\mu(A)=\min\{\lambda_{\max}(A),|\lambda_{\min}(A)|\}$.  Applying
hyperplane rounding to $I+sA$ and to $I-sA$ at their positive-semidefinite
spectral endpoints gives
\begin{equation}\label{eq:spectral-edge}
 M(A)\ge\frac2\pi n(n-1)\arcsin\!\left(\frac1{\mu(A)}\right).
\end{equation}
There is no sign cancellation because
$A_{ij}\arcsin(sA_{ij})=\arcsin s$.  The bound is exact for $A=J-I$.

\section{Upper bound and conference matrices}

For $q\equiv1\pmod4$ a prime power, the symmetric Paley conference matrix
$C$ of order $q+1$ satisfies
\[
 C^2=qI,
 \qquad \|C\|_2=\sqrt q,
 \qquad M(C)\le(q+1)\sqrt q.
\]
Primes congruent to $1$ modulo $4$ have relative gaps $o(x)$; the fixed
modulus statement follows, for example, from Thorner--Zaman
\cite{thorner-zaman}.  Restricting a nearby Paley matrix to a principal
$n\times n$ block and applying Lemma~\ref{lem:fill-in} proves the upper half
of Theorem~\ref{thm:window}.

Conference matrices are spectrally flat, but their Boolean maxima are not
uniformly close to the hyperplane-rounding value.  An internally audited
local-field calculation gives the following stronger statement.

\begin{theorem}[conference local improvement]\label{thm:conference-local}
For every sequence of symmetric conference matrices,
\[
 \liminf_n\frac{M(C_n)}{n\sqrt{n-1}}
 \ge \gamma_{\mathrm{loc}}
 =0.6894707435\ldots>\frac2\pi.
\]
If $A_n$ differs from $C_n$ in $h_n=o(n^{3/2})$ unordered edges, the same
coefficient holds for $A_n$.
\end{theorem}

The stability assertion follows from the exact Lipschitz estimate
\[
 |M(A)-M(C)|\le4h.
\]
For the local improvement, puncture the normalized conference matrix in the
all-ones direction, hyperplane-round, and use the limiting aligned-local-field
law
\[
 L=\sqrt{\frac2\pi}|G_0|+\sqrt{1-\frac2\pi}\,G_1.
\]
Independently flipping a negatively aligned site with probability
$\min((-L)_+/2,1)$ gives the stated gain.  The traffic input is of the type
developed in \cite{pesenti}; the conclusion above is nonperturbative at fixed
algorithmic depth.

There is also a coefficientwise conference free-energy theorem.  If
$H=Q_C/\sqrt{n-1}$, then for every fixed $r$,
\[
 \frac{\kappa_{2r}(H)}n
 \longrightarrow
 (-1)^{r-1}\frac{(2r)!}{4r}\operatorname{Cat}_{r-1}.
\]
These are the Taylor coefficients of the formal orthogonal-model function
\[
 G(\beta)=\frac14\left(
 \sqrt{1+4\beta^2}-1-
 \log\frac{1+\sqrt{1+4\beta^2}}2
 \right).
\]
This is \emph{not} a uniform deterministic free-energy comparison: the
Taylor series has radius $1/2$, while an improvement over
$\gamma_{\mathrm{loc}}$ would require large $\beta$.  At degree ten the first
class-sensitive invariant appears, but its normalized fixed-order
contribution is $O(1/n)$; fixed-order universality survives.

\section{Structure and arithmetic}

\begin{proposition}[block control]\label{prop:block-control}
For $[n]=S\sqcup T$,
\[
 \Cut(A_S)+\Cut(A_T)\le\Cut(A)
 \le\Cut(A_S)+\Cut(A_T)+\|A_{S,T}\|_{\infty\to1}.
\]
\end{proposition}

The lower inequality combines two block optimizers and chooses a simultaneous
sign flip so that the cross term is nonnegative.  The upper inequality
decomposes a global optimal signed cut into its internal and cross parts.

Principal restriction and arbitrary one-vertex extension give
\[
 F(n)\le F(n+1)\le F(n)+n.
\]
Thus $|L_{n+1}-L_n|=O(n^{-1/2})$, and the set of subsequential limits of
$(L_n)$ is a closed interval.  In particular, $H(n)=F(n)^{2/3}$ is
nondecreasing.  These facts alone do not imply convergence.

\subsection{Two-block composition and localization}

For $A\in\mathcal S_m$, write
\[
 f(A)=\max_x|Q_A(x)|,\qquad
 p(A)=\max_xQ_A(x),\qquad r(A)=-\min_xQ_A(x),
\]
and put $\Delta(A)=|p(A)-r(A)|$, so that
$f(A)=\max\{p(A),r(A)\}$.  Given $B\in\mathcal S_n$,
$W\in\{\pm1\}^{m\times n}$, and $\sigma\in\{\pm1\}$, form
\[
 P_{\sigma,W}=\begin{pmatrix}A&W\\W^{\mathsf T}&\sigma B\end{pmatrix}.
\]
The two configurations $(x,y)$ and $(x,-y)$ give opposite bridge fields.
Consequently the following identity is exact:
\begin{equation}\label{eq:exact-block-cap}
 f(P_{\sigma,W})=
 \max_{x,y}\bigl(
 |Q_A(x)+\sigma Q_B(y)|+|x^{\mathsf T}Wy|
 \bigr).
\end{equation}
Indeed, this is just
$\max\{|h+b|,|h-b|\}=|h|+|b|$ for each projective pair of
states.  In the bridge-free internal problem, choosing the relative
orientation gives the further exact formula
\begin{equation}\label{eq:orientation-saving}
 \min_{\sigma=\pm1}\max_{x,y}|Q_A(x)+\sigma Q_B(y)|
 =f(A)+f(B)-\min\{\Delta(A),\Delta(B)\}.
\end{equation}
For $\sigma=1$ the left-hand cap before minimization is
$\max\{p(A)+p(B),r(A)+r(B)\}$; for $\sigma=-1$ it is
$\max\{p(A)+r(B),r(A)+p(B)\}$.  Comparing these four sums proves
\eqref{eq:orientation-saving}, including tied extrema.
This saving concerns the internal direct sum only; adding a bridge still
requires control of the correlated maximum in \eqref{eq:exact-block-cap}.

Put
\[
 R(m,n)=\min_{W\in\{\pm1\}^{m\times n}}
 \|W\|_{\infty\to1},
 \qquad \mu_s=\E|\varepsilon_1+\cdots+\varepsilon_s|,
\]
where the $\varepsilon_i$ are independent uniform signs.

\begin{theorem}[rectangular bridge bounds]\label{thm:rectangular-bridge}
For all $m,n\ge2$,
\begin{align*}
 \max\{m\mu_n,n\mu_m\}\le R(m,n)
 \le\min\bigl\{&\sqrt{2(\log2)mn(m+n-1)},\\
 &m\mu_n+\sqrt{2mn(n-1)\log2},\\
 &n\mu_m+\sqrt{2mn(m-1)\log2}\bigr\}.
\end{align*}
In particular, if $n=o(m)$, then
\[
 R(m,n)=m\mu_n\bigl(1+O(\sqrt{n/m})\bigr).
\]
\end{theorem}

\begin{proof}
For fixed $W$,
$\|W\|_{\infty\to1}=\max_y\sum_i|(Wy)_i|$.  Averaging over
$y$ gives the lower bound $m\mu_n$; transposition gives $n\mu_m$.

For the first upper bound, take $W$ uniformly at random.  For each rank-one
sign pattern $(x_i y_j)_{ij}$, Hoeffding's inequality bounds the upper tail of
$x^{\mathsf T}Wy$.  There are $2^{m+n-1}$ distinct patterns, so a union bound
gives the first displayed estimate.  For the asymmetric estimate, fix $y$
and put $Z_y=\sum_i|(Wy)_i|$.  It has mean $m\mu_n$, and changing one entry
of $W$ changes $Z_y$ by at most $2$.  McDiarmid's inequality, followed by a
union bound over the $2^{n-1}$ projective choices of $y$, gives the second
upper bound; transposition gives the third.  At an equality threshold one
applies the argument with an arbitrarily small positive increment and then
uses finiteness.  Finally, the Khintchine lower bound
$\mu_n\gg\sqrt n$ compares the asymmetric error to $m\mu_n$.
\end{proof}

\begin{corollary}[unbalanced composition]\label{cor:unbalanced-composition}
Let $H(n)=F(n)^{2/3}$ and $N=m+n$, where $m\ge n\ge2$.  Then
\begin{equation}\label{eq:unbalanced-composition}
 H(N)\le H(m)+H(n)+1.291\sqrt{nN}.
\end{equation}
Thus, for every fixed $\varepsilon>0$, the condition
$n\le N^{1-\varepsilon}$ gives an $O(N^{1-\varepsilon/2})$ defect; merely
$n=o(N)$ gives an $o(N)$ defect.
\end{corollary}

\begin{proof}
Choosing optimal children and a bridge attaining $R(m,n)$ in
\eqref{eq:exact-block-cap} gives
$F(N)\le F(m)+F(n)+R(m,n)$.  Concavity, subadditivity of the
$2/3$ power, and $F(m)\ge m\sqrt{m-1}/\pi$ from
\eqref{eq:squared-gram-bound} yield
\[
 H(N)\le H(m)+H(n)
 +\frac23\pi^{1/3}\sqrt{2\log2}
 \left(\frac m{m-1}\right)^{1/6}\sqrt{n(N-1)}.
\]
The coefficient is at most
$(2/3)\pi^{1/3}\sqrt{2\log2}\,2^{1/6}
=1.2904012036\ldots<1.291$.
\end{proof}

The all-splits formulation of power-saving composition is stronger than
needed.  The following localization is the main structural gain.

\begin{theorem}[comparable-split localization]\label{thm:comparable-splits}
Let $(b_n)$ be a nonnegative nondecreasing sequence.  Suppose there are fixed
$C<\infty$, $\delta>0$, and $n_0$ such that
\[
 b_{r+s}\le b_r+b_s+C(r+s)^{1-\delta}
\]
whenever $r,s\ge n_0$ and $1/2\le r/s\le2$.  Then $b_n/n$ converges.
Consequently, this hypothesis only for comparable splits of
$b_n=H(n)=F(n)^{2/3}$ would prove convergence of $L_n$.
\end{theorem}

\begin{proof}
Fix $k\ge n_0$ and construct a block of size $qk$ from $q$ leaves of size
$k$, recursively splitting each leaf count into its floor and ceiling halves.
Every merge is between comparable children.  An internal node with $\ell$
descendant leaves contributes at most $C(\ell k)^{1-\delta}$.  Charge this
equally to those leaves.  Along any leaf-to-root path, the first internal size
is at least $2$ and subsequent sizes grow by a factor at least $3/2$.
Therefore each leaf receives at most
\[
 Ck^{1-\delta}K_\delta,
 \qquad
 K_\delta=\frac{2^{-\delta}}{1-(2/3)^\delta}.
\]
Summing over the leaves gives
\[
 \frac{b_{qk}}{qk}\le\frac{b_k}{k}+CK_\delta k^{-\delta}.
\]
For arbitrary $n$, take $q=\lceil n/k\rceil$ and use monotonicity
$b_n\le b_{qk}$.  Letting $n\to\infty$ gives the same upper bound for the
limsup of $b_n/n$.  Finally let $k\to\infty$ along a subsequence attaining
the liminf.  The limsup and liminf coincide.
\end{proof}

The compact finite audit emphasizes why a positive defect is unavoidable.
Literal subadditivity of $H$ fails already at $2+2$:
$H(4)-2H(2)=4^{2/3}-2=0.519842\ldots$.  With both children of order at
least $3$, it first fails at $3+5$ by
\[
 10^{2/3}-3^{2/3}-4^{2/3}=0.0416629\ldots>0.
\]
Among the $42$ pairs with $2\le m\le n$ and $m+n\le15$, using the
evidence-labelled exact-value table, $21$ have positive defect, including
$15$ of the $24$ comparable pairs.  The largest is
$H(11)-H(5)-H(6)=1.167629\ldots$.  Exhaustive bridge enumeration for four
fixed pairs of optimal children also finds orientation sensitivity at
$3+3$: the two relative orientations have exact optimized caps $7$ and $5$.
These are finite checks, not evidence for an asymptotic exponent; the keys,
integer sign certificates, and all $2^{mn-1}$ bridge searches are reproduced
by \texttt{experiments/composition\_small\_orders.py}.

The quadratic form satisfies
\[
 x^{\mathsf T}Ax\equiv n(n-1)\pmod4,
 \qquad
 F(n)\equiv\binom n2\pmod2.
\]
For odd $n$, every rectangular cut sum contains an even number of terms, so
$\Cut(A)$ is even.

\begin{theorem}[forced unbalancedness]
If $n\equiv3\pmod4$, then
\[
 F(n)>\min_{A\in\Sn}\Cut(A),
\]
and no $M$-minimizing matrix is balanced, where balanced means
$M_+(A)=M_-(A)$.
\end{theorem}

This explains all computed failures of $F(n)=\min_A\Cut(A)$ at
$n=3,7,11$ and forces another failure at $n=15$, independently of the
enumeration.

\section{Exact computation and provenance}\label{sec:computation}

Switching allows the first row and column to be normalized to $+1$ off the
diagonal.  The public iterator then enumerates
\[
 2^{(n-1)(n-2)/2}
\]
labeled switching-normalized matrices.  For each matrix, global-sign symmetry
reduces the Boolean cube to $2^{n-1}$ vectors.  This direct method reproduces
orders through $6$ and is intentionally simple.

The historical computations through order $15$ used a more efficient exact
integer pipeline: Gray-code energy updates, principal-submatrix pruning from
Lemma~\ref{lem:fill-in}, globally canonical switching/permutation
certificates, chunked extension levels, and threshold emptiness combined with
the congruence above.  The order-$15$ value $M=54$ was certified by two
independent empty searches at thresholds $46$ and $50$, overlapping
certificate sets, and direct evaluation of the witness
$[C_{14},\mathbf1]$ on the full antipodal half-cube.

Two implementation failures were discovered and bounded during that work.
An early 64-bit certificate encoding overflowed once
$\binom{k-1}{2}>64$; affected levels were recomputed using 128-bit
certificates.  A fixed-size resume buffer could overflow for sufficiently
large seed files; every certified run used smaller per-process seeds, and a
patched resume reproduced the relevant levels.  These incidents are recorded
because they materially affect how the exhaustive claims should be assessed.

The public repository does not ship the multi-gigabyte intermediate levels.
Accordingly, entries without a public external reproduction retain a
historical-audit status in \path{results/reported_exact_values.json}; the
evidence label on every entry should be read separately.

There is now independent public evidence for part of this range.  White's
order-$11$ package \cite{white-k11}, pinned at the revision recorded in
\texttt{results/external\_verification.json}, contains an explicit witness and
a standalone C++ lower-bound enumerator.  The present audit re-evaluated all
$2^{10}$ projective spin vectors for the witness and reran the lower search:
it visited $50{,}778{,}686$ nodes, completed $936{,}720$ candidates, and found
no signing with pair-sum cap $15$.  This independently proves $F(11)=17$.
An external extension witness then gives $F(12)\le18$; monotonicity, the
order-$11$ result, and even parity give equality.  Six symmetry-complete
CP-SAT cases excluding cap $18$ at order $13$ were also rerun from source,
and the order-$13$ and order-$14$ witnesses were exhaustively evaluated,
corroborating $F(13)=20$ and $F(14)=21$ \cite{beaumont-ledger}.  Those CP-SAT
runs do not emit standalone proof objects, so their evidence label remains
solver-certified rather than formal.  None of these external artifacts
replaces the historical order-$15$ threshold audit.

Heuristic search is not a certificate: at order $15$, many annealing restarts
stalled at $M=58$ although the exact optimum is $54$.

\section{The Hermite residual and the remaining absolute-field problem}
\label{sec:hermite}

The squared-Gram proof can be decomposed beyond its first Hermite chaos.  Let
$y_\sigma$ be the rounded sign vector for $X_\sigma$, write
\[
 y_\sigma=\sqrt{\frac2\pi}\,z_\sigma+\eta_\sigma,
 \qquad R_\sigma=\E[\eta_\sigma\eta_\sigma^{\mathsf T}],
\]
and couple the two orientations using one Gaussian vector $g$.  The arcsine
Schur expansion gives, for $n\ge5$,
\[
 \frac{R_++R_-}{2}\succeq\frac{1-2/\pi}{2}I.
\]
The paired mean local-field square is at least one, with equality only when
$B^2=I$.

Put
\[
 u=\frac{y_++y_-}{2},\qquad v=\frac{y_+-y_-}{2}.
\]
Coordinatewise exactly one of $u_i,v_i$ is nonzero:
\[
 (u_i,v_i)=
 \begin{cases}
 (\sgn g_i,0),&|g_i|>|(Bg)_i|,\\
 (0,\sgn(Bg)_i),&|(Bg)_i|>|g_i|.
 \end{cases}
\]
Define the complementary absolute field
\[
 U_n=\E\sum_i\max\{|(Bu)_i|,|(Bv)_i|\}.
\]
The exact remaining target is
\begin{equation}\label{eq:absolute-field}
 \liminf_n\frac{U_n}{n}\ge\sqrt{\frac2\pi}.
\end{equation}
A self-normalized cube-flip lemma converts
\eqref{eq:absolute-field} into a universal lower-bound improvement to
$0.3233368289\ldots$ for $L_n$, below the field-plus-spin constant now proved
above.  Separate second-moment floors
$\E\|Bu\|_2^2,\E\|Bv\|_2^2\ge n/2$ do not by themselves imply the required
first absolute moment.

Although \eqref{eq:absolute-field} is open universally, it is proved in a
large collection of regimes.  Write $D=B^2$ and
\[
 \overline\Delta_n=\frac{1}{n}\sum_{j\ne k}
 \max\{(n-1)^{-1/2},|D_{jk}|\}^3.
\]
The following statements are classified as \statusproved within the
project's internal audit.

\begin{enumerate}[label=(\roman*)]
\item If $\overline\Delta_n\to0$, then
      \eqref{eq:absolute-field} holds.  The same remains true after deleting
      $o(\sqrt n)$ exceptional source indices when the complementary cubic
      mass vanishes.
\item Fixed $k$-fold twin blow-ups of a globally decorrelated base satisfy
      the stronger coefficient
      $\sqrt{2/\pi}\sqrt{(k+2)/3}$.  Growing twin fibres are farther from
      extremality by a direct Boolean lift.
\item Disjoint strong matched pairs, under block-level global decorrelation,
      satisfy \eqref{eq:absolute-field}.
\item Correlation components of maximum size $265$ satisfy
      \eqref{eq:absolute-field} under block decorrelation.
\item A component partition is unnecessary: a strong-correlation graph of
      maximum degree $264$, with weak off-graph cubic mass, suffices.
\item More generally, with
      \[
      \mathcal L(\rho)=\frac2\pi
      \left(\arcsin\frac\rho2-\frac\rho2\right),
      \]
      the maximum-degree condition may be replaced by the oriented-edge
      density bound
      \[
      \frac1n\sum_{(j,k)\in E^{\mathrm{or}}}
      \frac{|\mathcal L(D_{jk})|}{|D_{jk}|}
      \le3.9698962518\ldots+o(1).
      \]
\end{enumerate}

The number $265$ is a method threshold, not a conjectured phase transition.
The unresolved regime must retain total cubic defect mass $\Omega(n)$ and a
dense or heavily overlapping strong-correlation graph after every admissible
weak-cubic decomposition.  Simple concentrated candidates such as twins,
matchings, bounded components, and bounded-degree giant components have all
been eliminated as obstructions.

\section{Barriers and corrected false routes}

The following negative conclusions are part of the research output.  They
prevent a successor from spending effort on arguments whose quantitative
ceiling is already understood.

\subsection{Spectral and degree-two certificates}

Every $A\in\Sn$ has $\|A\|_2\ge\sqrt{n-1}$.  Every diagonal $Y$ satisfying
$-Y\preceq A\preceq Y$ has
$\operatorname{tr}Y\ge n\sqrt{n-1}$.  For the pair of one-sided semidefinite
relaxations,
\[
 \SDP(A)\SDP(-A)\ge n^2(n-1),
\]
so $\max\{\SDP(A),\SDP(-A)\}\ge n\sqrt{n-1}$, with equality in the maximum
floor exactly for conference matrices.  Thus spectral, diagonal-scaling, and
degree-two SOS certificates cannot certify the desired constant below the
spectral ceiling.  The symmetric Grothendieck framework gives the relevant
norm dictionary \cite{friedland-lim}.

The tempting one-sided claim $\SDP(A)\ge n\sqrt{n-1}$ is false.  For the
order-$5$ matrix with one positive off-diagonal edge and all other edges
negative, exact primal and dual certificates give
\[
 \SDP(A)=9,
\]
not the earlier numerical estimate $9.008$.

\subsection{Composition, nesting, and regularity}

Separate PSD-Gram control of bilinear or Kronecker cross terms loses too much
to beat the spectral ceiling.  The admissible doubling construction also has
ratios that increase under iteration.  A generic expression
$A\otimes B+I\otimes S$ need not even be Seidel: some off-diagonal entries are
zero.

For a one-vertex extension $[A,a]$,
\[
 M([A,a])=max_y\left(|y^{\mathsf T}Ay|+2|a\cdot y|\right).
\]
If the maximizers of $A$ are isotropic, every extension costs at least
$2\sqrt n$ asymptotically, forcing a nested chain to coefficient at least
$4/3$.  This applies to several prominent conference-derived seeds.

Switching-regularity alone does not imply spectral attainment.  It gives a
$\{\pm1\}$ eigenvector, but the corresponding eigenvalue need not be
extremal.  The order-$9$ rook-graph Seidel matrix is regular with exact
$M=24<27=9\|A\|_2$.

\subsection{Single-split and superadditive methods}

For a rectangular sign matrix, the fractional relaxation of
$\|B\|_{\infty\to1}$ equals $q\E|S_p|$ exactly.  Bowlin's theorem is the
classical maximum-frustration form of this statement, including its equality
and uniqueness structure \cite{bowlin}.  Therefore improvement beyond the
first moment is a genuine integrality-gap problem, not an averaging or linear
duality problem.  The rigorous universal ceiling of the single-split route is
$2/(3\sqrt3)=0.384900\ldots$; the often quoted $0.335$ is only numerical.

Cut superadditivity and the associated free energy do not imply convergence.
An explicit analytic scalar sequence satisfies the relevant monotonicity,
increment, window, and superadditivity properties while oscillating.  Fekete
normalizations fail on exact small values, and bounded-difference
concentration has speed $n$ where overcoming the number of sign matrices
requires speed $n^2$.

\subsection{Several splits and Fourier summaries}

The several-splits functional
\[
 V(A)=\E_u\max_S
 \sum_{j\notin S}\left|\sum_{i\in S}A_{ij}u_i\right|
\]
lies between the single-split mean and $\Cut(A)$ and is strictly better at
several computed orders.  It has an exponential $L^2$ packing, but the
desired growing moments do not follow from low-degree Fourier mass.

Several explicit counterexamples delimit this route:
\begin{itemize}
\item for every fixed even swap distance $h$, there are Seidel matrices for
      which a split-field increment vanishes identically although the field
      itself is nonconstant;
\item dense, subset-robust level-two Fourier mass can coexist with constant
      modulus, so all $L^p$ norms are equal;
\item positive shifted absolute kernels have a level-four multiplier whose
      sign changes at the natural shift scale, ruling out a common
      level-four Gram sign;
\item low-degree mass, coefficient spread, and restriction robustness alone
      cannot force $\sqrt p$ moment growth.
\end{itemize}
Any successful continuation must use realizability of the actual shifted
absolute-linear-form process or average rigidity over a proportional packed
family, not only generic Fourier summaries.

\subsection{Other closed shortcuts}

Graphon convergence is blind because every candidate at the relevant scale
converges to the zero graphon.  Guerra-style interpolation has no disorder to
integrate and convexification admits the zero matrix.  Conference local-field
arithmetic gives exact finite ceilings but no uniform asymptotic gap; the
compatibility equation $C(Cx)=(n-1)x$ remains essential.  Fixed-order
conference cumulants agree with the orthogonal model, but that coefficientwise
fact does not justify a low-temperature comparison.

\section{Related formulations}

For the induced-subset discrepancy $D(n)$,
\[
 D(n)\le F(n)\le5D(n).
\]
Thus the Erd\H{o}s--Spencer problem has the same order but not the same
constant \cite{erdos-spencer}.

The coding formulation is exact.  Put $N=\binom n2$ and define the augmented
cut code
\[
 \mathcal C_n=
 \{(t+z_i+z_j)_{i<j}:t,z_1,\ldots,z_n\in\mathbb F_2\}
 \subseteq\mathbb F_2^N.
\]
If $\rho(\mathcal C_n)$ is its Hamming covering radius, then
\begin{equation}\label{eq:covering-radius}
 F(n)=N-2\rho(\mathcal C_n).
\end{equation}
Indeed, encode a signing by $a_{ij}=(-1)^{b_{ij}}$ and a spin vector by
$x_i=(-1)^{z_i}$.  Its pair sum is
$N-2\operatorname{wt}(b+(z_i+z_j)_{i<j})$; adjoining the complement bit $t$
accounts for the absolute value, and minimizing over $b$ takes the covering
radius.  This is the cut-code connection used in the order-$11$ certificate
\cite{white-k11}.  In particular $F(11)=17$ is equivalent to
$\rho(\mathcal C_{11})=19$, contradicting the covering-radius formula
conjectured by Esmaeili and Zaghian \cite{esmaeili-zaghian}.  General
weightwise nonlinearity is discussed by Gini and M\'{e}aux
\cite{gini-meaux}, but no exact weight-two asymptotic constant is known.

For $B\in\{\pm1\}^{p\times q}$,
\[
 \min_B\|B\|_{\infty\to1}
 =pq-2F_{\max}(K_{p,q}),
\]
where the right side is maximum frustration of a signed complete bipartite
graph.  Bowlin gives exact formulas for $p\le7$ \cite{bowlin}; the proportional
regime remains open.  This is closely related to the Gale--Berlekamp
switching game \cite{fishburn-sloane,carlson-stolarski,pellegrino-raposo}, but
that game is bilinear and its constants do not directly transfer to the
symmetric quadratic problem.

\section{Open problems and continuation map}

\begin{problem}[limit]
Does $\lim_n L_n$ exist?  If so, what is its value?
\end{problem}

\begin{problem}[power-saving composition]\label{prob:composition}
Put $H(n)=F(n)^{2/3}$.  Prove that there are fixed
$C<\infty$, $\delta>0$, and $n_0$ such that
\[
 H(m+n)\le H(m)+H(n)+C(m+n)^{1-\delta}
\]
whenever $m,n\ge n_0$ and $1/2\le m/n\le2$.  Theorem~\ref{thm:comparable-splits}
shows that comparable splits alone make $H(n)/n$ converge, and hence make
$L_n=(H(n)/n)^{3/2}$ converge.  Corollary~\ref{cor:unbalanced-composition}
already supplies a power saving when the smaller block is polynomially
smaller, but it gives only an order-$m+n$ error on comparable blocks.  External
composition campaigns identify the same target and reduce one
finite-temperature version to controlling a least-likely-output path in a
noisy rank-one code \cite{beaumont-ledger}.  The unresolved obstacle is
optimizer-sensitive: the internal block energies and the bridge field are
evaluated on the same spin assignment.
\end{problem}

\begin{problem}[absolute field]
Prove \eqref{eq:absolute-field} for the remaining dense overlapping
correlation regime, or construct a counterexample.  A proof would sharpen the
squared-Gram mechanism to $0.3233368289\ldots$, which is now below the
field-plus-spin constant $\gamma_*$; the problem remains structurally useful
but is no longer the leading universal-bound route.
\end{problem}

\begin{problem}[infinite sub-spectral family]
Is there a $\delta>0$ and an infinite sequence $A_n\in\Sn$ such that
$M(A_n)\le(1-\delta)n^{3/2}$?
\end{problem}

\begin{problem}[rectangular integrality gap]
Determine the proportional asymptotic of
$\min_B\|B\|_{\infty\to1}/(q\sqrt p)$.  The first concrete finite case beyond
Bowlin's formulas is $p=8$, in particular the $8\times16$ problem.
\end{problem}

\begin{problem}[order-15 cut]
Is $\min_{A\in\mathcal S_{15}}\Cut(A)$ equal to $24$ or $26$?  A structural
argument is preferable to repeating the infeasible threshold-$96$
enumeration.  A value-$24$ witness would have all fifteen order-$14$
deletions highly unbalanced.
\end{problem}

\begin{problem}[Bernoulli lower-tail rate]
Does
\[
 \lim_{n\to\infty}n^{-2}\log_2
 \#\{A\in\Sn:M(A)\le\lambda n^{3/2}\}
\]
exist?  Gaussian or typical-scale universality cannot by itself determine
this Bernoulli support-endpoint problem.
\end{problem}

\begin{problem}[growing-order conference comparison]
Can the fixed-cumulant conference universality be upgraded to a deterministic
comparison uniform in a growing order or inverse temperature?  Any such
argument must control the class-sensitive irreducible cores rather than
discard them termwise.
\end{problem}

\section{Reproduction guide}

The compact public computation is intentionally modest:
\begin{verbatim}
python -m pip install -e .
python -m unittest discover -s tests -v
python experiments/exact_small_n.py --max-order 6 \
  --output results/exact_small_n.json
python experiments/composition_small_orders.py \
  --output results/composition_small_orders.json
\end{verbatim}
It validates Seidel inputs, evaluates $M(A)$ on the antipodal half-cube,
enumerates every first-row-normalized labeled switching representative,
computes $\min_A M(A)$ exactly through order $6$, audits all scalar
composition defects through total order $15$, and exhausts the bridges for
four selected fixed-child cases.

The file \texttt{results/external\_verification.json} records the two pinned
external revisions used in the 2026-08-02 audit, SHA-256 hashes of the relevant
sources, exact commands, solver versions, and concise reproduced outputs.
The external repositories are intentionally not copied into this release.

The compact report source and the exact-result JSON are intended to be stable
handoff points.  Large future computations should publish a certificate
format, source revision, deterministic regeneration command, hashes, and an
independent evaluator before publishing a claimed optimum.

\section{Final status}

The limit problem is unresolved.  The strongest universal information is the
window $[\gamma_*,1/2]$, with
$\gamma_*=0.336493364432\ldots$, and exact values through order $15$.
The field-plus-spin theorem supersedes $1/\pi$ asymptotically, while the
squared-Gram result remains the stronger explicit finite statement of its
type.  Independent public computations now reproduce order $11$ and
corroborate orders $12$ through $14$; the order-$15$ optimum still rests on
the historical threshold audit.  For convergence, the cleanest live target
is now a power-saving composition bound only for comparable blocks; the
localization theorem proves that no all-splits estimate is needed.  No
candidate pair of separated infinite
subsequences is known.

\bibliographystyle{plain}
\bibliography{references}

\end{document}
